\subsection{Cayley-Hamilton and Minimal Polynomials}

\begin{theorem} \label{thm:cayley-hamilton-matrix} \hlt{Cayley--Hamilton Theorem}\\
Let $n\geq1$ and let $A\in\operatorname{Mat}(n,n,\mathbb{C})$ have characteristic polynomial $p_A(z)=\det(zI_n-A)$. Then $p_A(A)=0$.
\begin{proof}
\begin{enumerate}[label=(\roman*)]
\setlength{\itemsep}{0pt}
\item \hlt{Triangular Reduction}: By Theorem \ref{thm:simultaneous-schur-triangularization-commuting-family}(i) for $\{A\}$, choose a unitary $U$ such that $T\defeq U^*AU=[t_{ij}]$ is upper triangular.\\
Then $A=UTU^*$ by unitarity of $U$. Set $\lambda_j\defeq t_{jj}$ for $j\in\{1,\ldots,n\}$.\\
Then $p_A=p_T=\prod\limits_{j=1}^{n}(z-\lambda_j)$ by Theorem \ref{thm:similarity-preserves-characteristic-polynomial}(i) and Proposition \ref{prop:upper-triangular-characteristic-polynomial}(i).
\item \hlt{Polynomial Transfer}: For $k\in\mathbb{Z}_{\geq0}$, set $P_k\defeq(A^k=UT^kU^*)$. The identity $P_0$ holds by unitarity of $U$.\\
If $P_k$ holds, then $A^{k+1}=A^kA=UT^kU^*UTU^*=UT^{k+1}U^*$, so $P_{k+1}$ holds. Hence $(\forall\,k\in\mathbb{Z}_{\geq0})\,P_k$ by Theorem \ref{thm:mathematical-induction}.\\
For $j\in\{1,\ldots,n\}$, set $Q_j\defeq\prod\limits_{i=1}^{j}(T-\lambda_iI_n)$. By Definition \ref{def:polynomial-evaluation-on-square-matrix}, part (i), and the identities $P_k$, $p_A(A)=Up_A(T)U^*=UQ_nU^*$.
\item \hlt{Column Induction}: For $j\in\{1,\ldots,n\}$, set $R_j\defeq((\forall\,\ell\in\{1,\ldots,j\})\,Q_j\mathbf{e}_{\ell}=\mathbf{0})$.\\
The predicate $R_1$ holds because the first column of $T-\lambda_1I_n$ is zero. Suppose $j\in\{1,\ldots,n-1\}$ and $R_j$ holds.\\
By upper triangularity and the zero $(j+1,j+1)$ entry, $(\forall\,\ell\in\{1,\ldots,j+1\})\,(T-\lambda_{j+1}I_n)\mathbf{e}_{\ell}\in\operatorname{span}(\mathbf{e}_1,\ldots,\mathbf{e}_j)$.\\
By $R_j$ and linearity of $Q_j$, $(\forall\,\ell\in\{1,\ldots,j+1\})\,Q_{j+1}\mathbf{e}_{\ell}=Q_j(T-\lambda_{j+1}I_n)\mathbf{e}_{\ell}=\mathbf{0}$, so $R_{j+1}$ holds.\\
Hence $(\forall\,j\in\{1,\ldots,n\})\,R_j$ by Theorem \ref{thm:mathematical-induction}.
\item \hlt{Conclusion}: By $R_n$, every column of $Q_n$ is zero, so $Q_n=0$. Therefore $p_A(A)=UQ_nU^*=0$ by part (ii).
\end{enumerate}
\hlt{Structural Commentary}: Unitary triangularization transfers Cayley--Hamilton to a column-killing induction.\\
\hlt{Key Insight}: The $j$th factor extends the zero leading-column block from $j-1$ to $j$.
\end{proof}
\end{theorem}

Figure \ref{fig:cayley-hamilton-triangular-killing} depicts the invariant $R_j$: each new factor extends the zero leading-column block by one column.

\begin{illustrationfigure}[H]
  \centering
  \begin{tikzpicture}[
  >=Stealth,
  every node/.style={font=\small},
  mat/.style={draw=gray!55, fill=gray!5, rounded corners=2pt, inner sep=6pt, align=center},
  phasebox/.style={draw=blue!60!black, fill=blue!7, rounded corners=2pt, inner sep=5pt, align=center},
  arr/.style={->, line width=0.75pt, blue!70!black}
]
  \node[mat] (T) at (0,0)
    {$T=\begin{bmatrix}
      \lambda_1&*&*\\
      0&\lambda_2&*\\
      0&0&\lambda_3
    \end{bmatrix}$};
  \node[phasebox] (p1) at (3.7,0.8) {$P_1=T-\lambda_1I$\\first column $0$};
  \node[phasebox] (p2) at (3.7,-0.8) {$P_2=P_1(T-\lambda_2I)$\\first two columns $0$};
  \node[mat] (pn) at (7.7,0) {$P_n=\prod_i(T-\lambda_iI)=0$};

  \draw[arr] (T.east) -- (p1.west);
  \draw[arr] (p1.south) -- (p2.north);
  \draw[arr] (p2.east) -- (pn.west);
\end{tikzpicture}

  \caption{For upper-triangular $T$, the factors $T-\lambda_iI_n$ make the first $j$ columns of $Q_j$ vanish successively.}
  \label{fig:cayley-hamilton-triangular-killing}
\end{illustrationfigure}

\begin{theorem} \label{thm:cayley-hamilton-operator} \hlt{Cayley--Hamilton Theorem for Operators}\\
Let $\mathcal{V}\neq\{\mathbf{0}\}$ be finite-dimensional over $\mathbb{C}$, let $T\in\LL(\mathcal{V})$, and let $q_T$ be the characteristic polynomial of $T$. Then $q_T(T)=0$.
\begin{proof}
Choose a basis $\beta$ of $\mathcal{V}$ and set $A=[T]_{\beta\leftarrow\beta}$. The identity $q_T=p_A$ holds by Definition \ref{def:operator-characteristic-polynomial}.\\
The identity $[q_T(T)]_{\beta\leftarrow\beta}=q_T(A)$ holds by Propositions \ref{prop:matrix-representation-respects-linear-combinations} and \ref{prop:composition-and-matrix-multiplication}. Hence $[q_T(T)]_{\beta\leftarrow\beta}=p_A(A)=0$ by Theorem \ref{thm:cayley-hamilton-matrix}.\\
Therefore $q_T(T)=0$ by Theorem \ref{thm:linear-maps-and-matrices-same-data}.\\
\hlt{Structural Commentary}: Matrix representation transfers operator polynomial evaluation to the matrix Cayley--Hamilton identity.\\
\hlt{Key Insight}: The coordinate matrix of $q_T(T)$ is exactly $q_T(A)$.
\end{proof}
\end{theorem}

\begin{corollary} \label{cor:polynomial-reduction-by-characteristic-polynomial} \hlt{Reduction Modulo the Characteristic Polynomial}\\
Let $n\geq1$, let $A\in\operatorname{Mat}(n,n,\mathbb{C})$, and let $f,q,r\in\mathbb{C}[z]$. If $f=p_Aq+r$ with $r=0$ or $\deg r<n$, then $f(A)=r(A)$.
\begin{proof}
By the polynomial identity in the hypothesis, $f(A)=p_A(A)q(A)+r(A)$. Since $p_A(A)=0$ by Theorem \ref{thm:cayley-hamilton-matrix}, $f(A)=r(A)$.\\
\hlt{Structural Commentary}: Evaluation transfers division by $p_A$ to equality of the corresponding matrix polynomials.\\
\hlt{Key Insight}: The characteristic-polynomial term vanishes at $A$.
\end{proof}
\end{corollary}

\begin{corollary} \label{cor:powers-of-matrix-span-dimension-at-most-n} \hlt{Powers Span at Most $n$ Dimensions}\\
Let $n\geq1$, let $A\in\operatorname{Mat}(n,n,\mathbb{C})$, and set $\mathcal{W}_A\defeq\operatorname{span}(I_n,A,\ldots,A^{n-1})$. Then:
\begin{enumerate}[label=(\roman*)]
\setlength{\itemsep}{0pt}
\item $\operatorname{span}(I_n,A,A^2,\ldots)=\mathcal{W}_A$;
\item $\dim(\mathcal{W}_A)\leq n$.
\end{enumerate}
\begin{proof}
\begin{enumerate}[label=(\roman*)]
\setlength{\itemsep}{0pt}
\item Write $p_A(z)=z^n+\sum\limits_{j=0}^{n-1}c_jz^j$ by Theorem \ref{thm:matrix-characteristic-polynomial-basic-coefficients}(i). Then $A^n=-\sum\limits_{j=0}^{n-1}c_jA^j\in\mathcal{W}_A$ by Theorem \ref{thm:cayley-hamilton-matrix}.\\
If $0\leq j<n-1$, then $AA^j=A^{j+1}\in\mathcal{W}_A$; the remaining generator satisfies $AA^{n-1}=A^n\in\mathcal{W}_A$.\\
Hence $A\mathcal{W}_A\subseteq\mathcal{W}_A$ by linearity. Set $P_k\defeq(A^k\in\mathcal{W}_A)$ for $k\in\mathbb{Z}_{\geq0}$.\\
The predicate $P_0$ holds by $A^0=I_n\in\mathcal{W}_A$, and $P_k\Rightarrow P_{k+1}$ by $A\mathcal{W}_A\subseteq\mathcal{W}_A$. Thus $(\forall\,k\in\mathbb{Z}_{\geq0})\,P_k$ by Theorem \ref{thm:mathematical-induction}.\\
Therefore $\operatorname{span}(I_n,A,A^2,\ldots)\subseteq\mathcal{W}_A$; the reverse inclusion holds by the definition of $\mathcal{W}_A$.
\item The list $(I_n,A,\ldots,A^{n-1})$ spans $\mathcal{W}_A$ and has length $n$. It contains a basis of $\mathcal{W}_A$ by Theorem \ref{thm:reducing-spanning-lists}, so $\dim(\mathcal{W}_A)\leq n$.
\end{enumerate}
\hlt{Structural Commentary}: The characteristic recurrence makes $\mathcal{W}_A$ invariant; a basis sublist supplies the dimension bound.\\
\hlt{Key Insight}: $A\mathcal{W}_A\subseteq\mathcal{W}_A$ propagates containment to every matrix power.
\end{proof}
\end{corollary}

\begin{corollary} \label{cor:inverse-from-characteristic-polynomial} \hlt{Inverse from the Characteristic Polynomial}\\
Let $n\geq1$, let $A\in\operatorname{Mat}(n,n,\mathbb{C})$, and write $p_A(z)=\sum\limits_{j=0}^{n}c_jz^j$ with $c_n=1$. If $A$ is invertible, then $A^{-1}=-c_0^{-1}\sum\limits_{j=1}^{n}c_jA^{j-1}$.
\begin{proof}
Since $A$ is invertible, $0\notin\operatorname{spec}(A)$ by Corollary \ref{cor:matrix-invertible-iff-zero-not-spectrum}. Hence $c_0=p_A(0)\neq0$ by Theorem \ref{thm:matrix-characteristic-polynomial-basic-coefficients}(ii).\\
By Theorem \ref{thm:cayley-hamilton-matrix}, $c_0I_n=-\sum\limits_{j=1}^{n}c_jA^j$. By left multiplication by $c_0^{-1}A^{-1}$, $A^{-1}=-c_0^{-1}\sum\limits_{j=1}^{n}c_jA^{j-1}$.\\
\hlt{Structural Commentary}: Cayley--Hamilton separates the constant term, while invertibility permits multiplication by $c_0^{-1}A^{-1}$.\\
\hlt{Key Insight}: The nonzero constant coefficient permits solving the characteristic identity for $A^{-1}$.
\end{proof}
\end{corollary}

\begin{theorem} \label{thm:minimal-polynomial-basic-properties} \hlt{Basic Properties of the Minimal Polynomial}\\
Let $n,d\geq1$, let $A\in\operatorname{Mat}(n,n,\mathbb{C})$, and let $\lambda_1,\ldots,\lambda_d$ be the distinct eigenvalues of $A$.
\begin{enumerate}[label=(\roman*)]
\setlength{\itemsep}{0pt}
\item $\mu_A\mid p_A$;
\item $\deg\mu_A\leq n$;
\item $(\exists\,q_1,\ldots,q_d\in\mathbb{Z}_{>0})[(\forall\,i\in\{1,\ldots,d\})\,q_i\leq m_A(\lambda_i)\ \land\ \mu_A(z)=\prod\limits_{i=1}^{d}(z-\lambda_i)^{q_i}]$.
\end{enumerate}
\begin{proof}
\begin{enumerate}[label=(\roman*)]
\setlength{\itemsep}{0pt}
\item The polynomial $p_A$ is nonzero by Theorem \ref{thm:matrix-characteristic-polynomial-basic-coefficients}(i), and $p_A(A)=0$ by Theorem \ref{thm:cayley-hamilton-matrix}.\\
Hence $\mu_A\mid p_A$ by Theorem \ref{thm:operator-annihilator-iff-multiple-minimal-polynomial} and Definition \ref{def:minimal-polynomial-matrix}.
\item By (i) and $\deg p_A=n$ from Theorem \ref{thm:matrix-characteristic-polynomial-basic-coefficients}(i), $\deg\mu_A\leq n$.
\item By Theorem \ref{thm:simultaneous-schur-triangularization-commuting-family}(i) for $\{A\}$, choose a unitary $U$ such that $R\defeq U^*AU$ is upper triangular.\\
By Theorem \ref{thm:similarity-preserves-characteristic-polynomial}(i),(ii) and Proposition \ref{prop:upper-triangular-characteristic-polynomial}(i),(iii), $p_A(z)=\prod\limits_{i=1}^{d}(z-\lambda_i)^{m_A(\lambda_i)}$.\\
By factorization over $\mathbb{C}$ and monicity of $\mu_A$, write $\mu_A$ as a product of monic linear factors.\\
By (i), each factor occurs in $p_A$ with no greater exponent, so $\mu_A(z)=\prod\limits_{i=1}^{d}(z-\lambda_i)^{q_i}$ for integers $0\leq q_i\leq m_A(\lambda_i)$.\\
Fix $i\in\{1,\ldots,d\}$ and choose $\mathbf{v}\neq\mathbf{0}$ satisfying $A\mathbf{v}=\lambda_i\mathbf{v}$ by the eigenvalue hypothesis.\\
Then $\mathbf{0}=\mu_A(A)\mathbf{v}=\mu_A(\lambda_i)\mathbf{v}$ by Proposition \ref{prop:polynomial-eigenvector-evaluation}, so $\mu_A(\lambda_i)=0$ and $q_i\geq1$.\\
Since $i$ was arbitrary, $(\forall\,i\in\{1,\ldots,d\})\,q_i\geq1$, which proves (iii).
\end{enumerate}
\hlt{Structural Commentary}: Cayley--Hamilton supplies divisibility; Schur form and eigenvectors determine the exponents.\\
\hlt{Key Insight}: Every eigenvalue must occur in $\mu_A$, but never above its multiplicity in $p_A$.
\end{proof}
\end{theorem}

Figure \ref{fig:minimal-polynomial-divisibility} depicts $\mu_A$ as a common divisor of $p_A$ and every polynomial $f$ satisfying $f(A)=0$.

\begin{illustrationfigure}[H]
  \centering
  \begin{tikzpicture}[
  >=Stealth,
  every node/.style={font=\small},
  box/.style={draw=gray!55, fill=gray!5, rounded corners=2pt, align=center, inner xsep=5pt, inner ysep=4.5pt, minimum width=1.55cm},
  main/.style={draw=blue!60!black, fill=blue!7, rounded corners=2pt, align=center, inner xsep=5pt, inner ysep=4.5pt, minimum width=1.55cm},
  arr/.style={->, line width=0.75pt, blue!70!black}
]
  \node[main] (mu) at (0,0) {$\mu_A$};
  \node[box] (char) at (2.60,0.55) {$p_A$};
  \node[box] (other) at (2.60,-0.55) {$f(A)=0$};

  \draw[arr] (mu) -- (char);
  \draw[arr] (mu) -- (other);
\end{tikzpicture}

  \caption{$\mu_A$ is the least monic annihilator; arrows indicate that $\mu_A$ divides $p_A$ and every annihilator $f$.}
  \label{fig:minimal-polynomial-divisibility}
\end{illustrationfigure}
